% FILE: 05_evidence_receipts_and_gates.tex
% Project: lifecycle
% Role: process-lifecycle-state-definitions-and-transitions

\section{Evidence, Receipts, and Gates}
\label{sec:lifecycle-evidence}

\subsection{Tests}
\label{sec:lifecycle-evidence-tests}

The lifecycle project defines four test assertion blueprints in
\texttt{checkpoint\_\_lifecycle\_model\_complete.md}, expressed as Python pseudocode.
These blueprints constitute a formal specification of what a conforming test suite must
verify, but they are not executable test files. No \texttt{.py}, \texttt{.rs}, or
test runner files exist within the \texttt{lifecycle/} directory. The gap between
specification and execution is a known open item, documented explicitly in the
checkpoint status and in the open questions section of this chapter.

The four assertion blueprints are: \texttt{assert\_soundness(N)}, which verifies that
the net is a valid WF-net, 1-bounded, live, and free of dead transitions;
\texttt{assert\_fitness(N, L, threshold=0.95)}, which computes alignment fitness and
asserts it meets the threshold; \texttt{assert\_no\_ghost\_transitions(N, L)}, which
verifies that every labeled transition is fired at least once in the log or is classified
as a silent transition; and \texttt{assert\_decommission\_receipt(R\_d, K\_\mathrm{pub})},
which re-computes the model hash and log hash and verifies the receipt signature.

The checkpoint verdict rule is stated as a formal conjunction:
\[
\mathrm{status}(N) = \texttt{ALIVE} \iff
  \texttt{assert\_soundness}(N) \land
  \texttt{assert\_fitness}(N, L, 0.95) \land
  \texttt{assert\_no\_ghost\_transitions}(N, L)
\]
This conjunction is a formal specification, not a passed test. The lifecycle project
does not claim ALIVE at the module level because no executable test suite has run these
assertions against real Petri nets and real event logs in a running \texttt{wasm4pm}
instance. The PARTIAL status reflects this distinction precisely.

The adversarial test requirements are specified in the Van der Aalst Constitution
(Chicago TDD doctrine): inject impossible logs and verify rejection; verify that
model-vs-log mismatch is refused, not warned; verify temporal conformance so that
stages occur in lawful order with no impossible overlaps. These requirements are
specified but not yet executed.

\subsection{Receipts}
\label{sec:lifecycle-evidence-receipts}

The LIFECYCLE\_RECEIPT is recorded in \texttt{receipts/RECEIPT\_REGISTRY.md} with the
following fields:

\begin{itemize}
  \item \textbf{name}: LIFECYCLE\_RECEIPT
  \item \textbf{produced\_by}: lifecycle/ phase definition workflow
  \item \textbf{witness}: process mining lifecycle literature + wasm4pm-compat typestate
  \item \textbf{result\_type}: 12+ lifecycle phases defined with compat/wasm4pm coverage
    per phase
  \item \textbf{criteria\_met}: lifecycle $\geq 8$ files (actual: 41+); every phase has
    name, input type, output type, wasm4pm-compat state tag, admission requirement;
    phases cover ingest, extract, clean, admit, discover, conform, predict, export,
    archive, diligence
  \item \textbf{gaps\_remaining}: none listed
\end{itemize}

Five additional JSON receipts exist in the parent \texttt{receipts/} directory that
are relevant to the lifecycle framework's M\&A claims:
\texttt{rec\_ebitda\_rework\_001.json} (EBITDA rework cost quantification),
\texttt{rec\_wc\_ar\_002.json} (working capital accounts receivable),
\texttt{rec\_risk\_sla\_003.json} (SLA risk quantification),
\texttt{rec\_risk\_compliance\_004.json} (compliance risk), and
\texttt{rec\_residual\_standard\_005.json} (process residual standard). These are
sample receipts; the RECEIPT\_REGISTRY notes that production receipts require live
operational data.

The BLAKE3 hash-chaining ledger $H_n = \mathrm{BLAKE3}(\mathrm{Payload}_n \parallel
H_{n-1})$ is the specification for how receipts accumulate provenance across the full
lifecycle. The receipt chain is tamper-evident: any modification to a past receipt
invalidates all subsequent receipts in the chain. The receipt chain specification is
defined in the lifecycle doctrine but the running ledger is a \texttt{wasm4pm}
execution obligation.

\subsection{Checkpoints}
\label{sec:lifecycle-evidence-checkpoints}

The program-level checkpoint \texttt{PROCESS\_INTELLIGENCE\_ALIVE\_001.md} issues the
ALIVE verdict at the research program level with verification status code \texttt{0x00}
(SUCCESS\_ALIVE). The checkpoint confirms the following mathematical invariants:
admissibility boundary $R \vdash P_i = \mu(O^*, T, L)$, autonomic actuation
$\alpha(K, P, L, T) \to \tau$, token game fitness
$f(L, N) = 1 - \frac{\sum m}{\sum c} - \frac{\sum r}{\sum p}$, OCPQ refinement
$p \sqsubseteq_L c$, and decommissioning retirement $\delta(P)$.

The lifecycle-specific checkpoint \texttt{checkpoint\_\_lifecycle\_model\_complete.md}
defines the module-level ALIVE criteria as the four-assertion conjunction described
above. This checkpoint does not carry an ALIVE verdict stamp; it specifies the criteria
that would warrant one. The distinction between the program-level ALIVE checkpoint and
the module-level criteria document is material: the program-level checkpoint is a
research program milestone; the module-level ALIVE requires executable test passage.

The \texttt{lifecycle/audit\_\_lifecycle\_completeness.md} provides the completeness
audit checklist, verifying that every stage and taxonomy is covered by at least one
file with the required fields populated.

\subsection{ALIVE/PARTIAL/BLOCKED Status}
\label{sec:lifecycle-evidence-status}

The lifecycle project's current status is \textbf{PARTIAL}.

The PARTIAL verdict reflects the following precise state. The doctrine surface is
complete: 43 files covering all 13 lifecycle stages, all six quality gates, all five
taxonomies, the MAPE-K mapping, the maturity model, and the receipt specification.
The LIFECYCLE\_RECEIPT is satisfied with no gaps remaining. The type definitions in
\texttt{wasm4pm-compat} provide the type-level enforcement of the evidence lifecycle
state machine. The program-level ALIVE checkpoint (\texttt{0x00}) is issued.

The PARTIAL verdict is warranted because the following criteria are not yet met. First,
no executable test suite has run \texttt{assert\_soundness}, \texttt{assert\_fitness},
\texttt{assert\_no\_ghost\_transitions}, or \texttt{assert\_decommission\_receipt}
against real Petri nets and real event logs. Second, the \texttt{wasm4pm} execution
engine---required for online token replay, alignment computation, variant enumeration,
repair plan evaluation, actuation, receipt emission, model storage, and metric
time-series storage---is a separate downstream product that has not yet closed the
MAPE-K loop at runtime. Third, the Board Projection Holography framework does not yet
have a formal connection to the receipt chain or \texttt{wasm4pm} types established
through a compiled, running bridge.

No BLOCKED conditions exist: the path to ALIVE is well-specified and the required
downstream work is documented in the \texttt{wasm4pm} obligations per lifecycle phase.
